\section*{摘要}
\addcontentsline{toc}{section}{摘要}
本文面向课程给出的众包离线高精度位置聚合开放问题，开发了使用实际大语言模型的科研智能体CrowdLocateLab。系统将问题阅读、可检验假设、组合方法注册、数值实验、失败诊断、独立审阅调用、数学检查和方法冻结串联起来，保留原始模型输出、工具输入输出及可恢复状态。两轮研究合计完成25个研究动作、27次实际模型调用和12次实验或消融调用；第二轮以第一轮失败为已知先验，再用新数据评价修订。

研究关注污染测距与历史观测下的当前位置估计，实现静态与恒速模型、距离与对数残差、稳健损失、截尾、子集初始化及观测交叉验证等可组合组件。理论部分给出自加权标量偏差、镜像不可辨识性、局部误差界，并推导镜像世界下接受率与错误率的限制。校准部分对比全局与可观测选模分组的分割共形区域，明确有限样本秩和保证条件。各数学结论的适用范围与实证结果分开陈述。

第一轮主测试显示总体改善，但静态退化、稀疏模型排除和子群欠覆盖促成第二轮修订。第二轮在720例新主测试上的平均误差为3.920米，相比WLS降低40.99\%；相对第一轮冻结方法的改善为0.145米，配对95\%区间$[-0.043,0.334]$跨零。8份观测的运动边界中，误差由10.788米降至6.217米，但仍高于恒速稳健基线的5.107米。静态正常误差相对WLS退化29.86\%，预设10\%容忍目标未达到。分组区域在1200例独立覆盖检查中，90\%与95\%名义水平的实际覆盖分别为88.17\%与95.00\%；分组没有消除真实运动异常场景的欠覆盖。

两轮冻结后评测及敏感性分析合计完成22052次定位计算，129项软件测试通过。交付包括完整源码、真实运行证据、逐例数据、只读研究浏览器、新观测定位入口及LaTeX源码。合成研究只回答原工业问题的一个可计算子问题，未完成完整工业题验收，也不将已有统计方法的组合称为原创通用定理。

\vspace{4mm}
\noindent\textbf{关键词：}科研智能体；异步定位；稳健估计；模型选择；可辨识性；分组共形预测
